\chapter{基于强化学习的多无人机协同运输控制方法}

单架无人机运输已在低空物流场景中得到广泛应用，适用于轻小载荷与中短距离任务，具有部署灵活、组织简洁的特点。当运输对象的重量和体积进一步增加时，单机系统在推力储备、续航能力和安全裕度上的限制会迅速显现，运输能力边界也随之暴露。超出单机承载范围的任务往往需要通过分拆载荷或多次往返完成，任务效率与执行稳定性均会受到明显影响。

多机协同负载运输为突破单机承载上限提供了可行途径。多架无人机经柔性绳索共同牵引同一负载后，系统的整体负载能力得到提升，控制问题也转化为多体耦合条件下的协同控制问题。推力分配、绳索张力和负载运动在运输过程中持续耦合，负载摆动、绳索松弛与拉紧切换、机间相对位形变化以及输入约束同时进入系统演化，使协同运输表现出欠驱动、强非线性和多体耦合特征。面向这一问题，本章提出基于强化学习的多无人机协同运输控制方法，并围绕三机---负载系统在绳索约束下的协同运输问题展开研究。章节内容依次覆盖耦合动力学建模、共享回报的多智能体序贯决策表述、局部观测与归一化控制量表示及奖励函数设计、多智能体近端策略优化（Multi-Agent Proximal Policy Optimization，MAPPO）形式的集中训练与去中心化执行框架，以及基于 MJX 的并行训练加速实现。实验部分采用与前述任务定义一致的参数共享近端策略优化（Proximal Policy Optimization，PPO）完成训练评估，并通过定点运输任务与圆周参考扩展任务验证该方法在复杂耦合系统稳定控制与任务适应性方面的有效性，整体控制框架如图~\ref{framework3} 所示。

\begin{figure}[!h]
  \centering
  \includegraphics[width=\textwidth]{pic/算法框图3.png}
  \caption{基于强化学习的多无人机协同运输控制方法框架图}
  \label{framework3}
\end{figure}

\section{问题描述}

三机协同吊运同时包含负载平动、无人机姿态演化以及绳索几何约束，系统状态、控制输入和约束条件需在统一坐标框架下加以描述。本节针对协同运输任务给出系统对象与建模假设，定义负载与绳索相关变量，建立系统动力学模型，并据此明确控制目标与约束条件。在三机协同运输过程中，负载运动并非由单架无人机独立决定，而是由三架无人机经绳索传递的联合受力共同塑造。单机推力变化会改变绳索方向与张力分布，负载受力状态随之调整，其余无人机的工作点也同步变化，后续观测构造、动作映射与奖励设计均建立在这一耦合关系之上。协同运输系统由三架四旋翼无人机和一个负载构成，三架无人机分别通过柔性绳索连接至同一负载，三根绳索均接于负载质心，负载在建模中按质点处理，研究对象限定为负载平动运输过程，即系统在满足飞行安全、绳索约束与输入约束的条件下，将负载稳定送入目标区域，并在目标附近实现速度衰减与摆动抑制。运输过程可划分为构型建立、负载推进与目标附近稳定保持三个阶段，运输初期需要形成稳定牵引构型并建立有效合力，中段运输要求三架无人机维持协调张力分布，目标附近则以速度衰减和残余摆动抑制为主要特征；任一无人机姿态或推力的变化都会改变绳索方向与张力，负载受力状态随之发生调整，其他无人机的控制需求也相应变化。三机协同运输的系统性能最终体现为整体运输效果，无人机推力与姿态调整通过绳索方向和张力变化传递至负载，负载状态再经绳索几何关系反馈至其余无人机，多机之间形成显著的间接耦合，共享目标与强耦合约束在该任务中同时存在，协同运输据此表现为围绕单一负载目标展开的多体控制问题。

多无人机协同吊运系统中的绳索建模通常存在两类典型处理方式。第一类将绳索视为无质量、不可伸长且始终张紧的刚性连杆约束，利用绳索方向单位向量和非负张力描述无人机与负载之间的力传递关系，该类模型在多四旋翼协同悬挂负载动力学与几何控制研究中得到了广泛采用\cite{sreenath2013dynamics,lee2013geometricCooperating}。第二类将柔性绳索离散为多段串联刚性杆，并通过球铰连接各段连杆，以刻画绳索在空间中的弯曲和大变形特征；Goodarzi和Lee将连接多无人机与刚体负载的柔性绳索表示为任意数量串联连杆系统\cite{goodarzi2015dynamicsFlexibleCables}，Kotaru等进一步将柔性绳索建模为带球铰的有限段连杆，并在此基础上分析多无人机—负载系统的微分平坦性与控制问题\cite{kotaru2018differentialFlatnessFlexibleCables}。本文面向三机协同运输控制问题，重点关注绳索张紧状态下的力传递和负载平动控制，因此采用绳索约束下的多机—负载耦合动力学模型。

为刻画协同运输系统的受力关系与耦合特性，进一步定义负载与绳索相关变量如下。负载视为质点，其位置与速度分别记为 $p_L\in\mathbb R^3$ 和 $v_L\in\mathbb R^3$；第 $i$ 架无人机的位置、速度、姿态矩阵和角速度分别记为 $p_i\in\mathbb R^3$、$v_i\in\mathbb R^3$、$R_i\in SO(3)$ 和 $\Omega_i\in\mathbb R^3$。第 $i$ 根绳索的向量、长度与单位方向定义为
\begin{equation}
    \ell_i = p_i - p_L,\qquad
    l_i = \|\ell_i\|,\qquad
    q_i = \frac{\ell_i}{\|\ell_i\|}.
\end{equation}
负载目标位置记为 $p_g\in\mathbb R^3$，相应的位置误差与速度误差写为
\begin{equation}
    e_p = p_L - p_g,\qquad
    e_v = v_L.
\end{equation}
据此，第 $i$ 架无人机状态、负载状态和系统整体状态分别记为
\begin{equation}
    x_i = (p_i,v_i,R_i,\Omega_i),\qquad
    x_L = (p_L,v_L),\qquad
    x = (x_1,x_2,x_3,x_L).
\end{equation}
绳索作为柔性连接件，仅在拉紧状态下传递拉力。负载并不直接受控，三架无人机通过绳索张力共同改变其受力状态。第 $i$ 根绳索张力记为 $\lambda_i\ge 0$，其方向沿 $q_i$ 指向无人机。为表征绳索的单边约束特征，第 $i$ 根绳索的张力建模为
\begin{equation}
    \lambda_i = \max\Big(0,\; k_i(l_i-l_{0,i}) + c_i\dot l_i\Big),
\end{equation}
其中 $l_{0,i}$ 为绳索自然长度，$k_i$ 为等效刚度系数，$c_i$ 为等效阻尼系数。
绳索伸长速度由无人机与负载沿绳索方向的相对运动决定，可表示为
\begin{equation}
    \dot l_i = q_i \cdot (v_i-v_L).
\end{equation}
绳索处于松弛状态时有 $\lambda_i=0$；沿绳索方向继续拉伸后，弹性项与阻尼项共同决定张力大小。三根绳索均连接于负载质心，负载动力学只保留平动部分。
第 $i$ 架无人机的平动动力学写为
\begin{equation}
    m_i\ddot p_i = f_i R_i \mathbf e_z - m_i g\mathbf e_z - \lambda_i q_i + d_i,
\end{equation}
其中 $m_i$ 为无人机质量，$f_i$ 为总推力，$d_i$ 为平动扰动项。其转动动力学写为
\begin{equation}
    J_i\dot \Omega_i + \Omega_i \times J_i\Omega_i = \tau_i + d_{\tau,i},
\end{equation}
其中 $J_i$ 为转动惯量矩阵，$\tau_i$ 为控制力矩，$d_{\tau,i}$ 为转动扰动项。
负载作为质点处理，其动力学仅包含平动部分，即
\begin{equation}
    m_L \ddot p_L = \sum_{i=1}^{3}\lambda_i q_i - m_L g\mathbf e_z + d_L,
\end{equation}
其中 $m_L$ 为负载质量，$d_L$ 为负载平动扰动项。
系统动力学中存在明确的受力传递链。无人机推力经姿态矩阵投影到世界坐标系后形成机体加速度，绳索张力再把单机动作耦合到负载运动。任一无人机的控制变化都会改变负载受力状态和其余无人机的工作点，协同运输由此表现出欠驱动多体系统的显著耦合特征。在该动力学框架下，负载运动由三架无人机的联合受力共同塑造，而单机对负载的作用能力又受推力约束、瞬时构型、绳索方向以及邻机动作共同影响，协同控制关注多机联合受力对负载运动的整体组织，单机误差调节只是这一协同过程中的局部表现。

基于上述任务场景与动力学描述，协同运输控制问题可进一步表述为负载到达与稳定保持的联合调节问题。具体而言，控制目标在于运输过程中保持较小的 $\|e_p\|$ 与 $\|e_v\|$，并抑制绳索摆动。为度量摆动程度，定义最大绳索摆角为
\begin{equation}
    \theta_{\max} = \max_i \arccos(q_i\cdot \mathbf e_z).
\end{equation}
除负载位置稳定性外，系统同时受到单机可实现性与运输安全约束。对第 $i$ 架无人机，控制输入满足
\begin{equation}
    f_i \in \mathcal F_i,\qquad
    \tau_i \in \mathcal T_i,\qquad
    \psi_i^d \in \Psi_i,
\end{equation}
其中 $\mathcal F_i$、$\mathcal T_i$ 和 $\Psi_i$ 分别表示总推力、姿态力矩和偏航指令的允许集合。无人机姿态约束写为
\begin{equation}
    |\phi_i| \le \Phi_{\max},\qquad
    |\theta_i| \le \Theta_{\max}.
\end{equation}
系统安全约束写为
\begin{equation}
    d_{ij}\ge d_{\min},\qquad
    z_i \ge z_{\min}^u,\qquad
    z_L \ge z_{\min}^L,\qquad
    \theta_{\max}\le \theta_{\mathrm{lim}},
\end{equation}
其中 $d_{ij}$ 表示任意两架无人机之间的相对距离，$z_i$ 和 $z_L$ 分别表示无人机与负载的高度。
在上述动力学与约束条件下，协同控制策略需要持续减小负载位置偏差、速度偏差与摆动幅度，并控制输入消耗。该目标与后续强化学习中的共享奖励优化目标保持一致，协同运输任务也据此被写成约束条件下的序贯决策问题。

\section{协同运输任务的强化学习建模}

三机协同运输中的控制输入由多架无人机分别产生，而运输结果由联合动作共同决定。执行阶段依赖局部可得信息，训练阶段则需要全局状态刻画多机耦合关系。本节在协同运输动力学基础上，将任务表述为去中心化部分可观测马尔可夫决策过程（Decentralized Partially Observable Markov Decision Process，Dec-POMDP），并给出与该任务相匹配的观测空间、控制量表示和奖励结构。相较于单机轨迹跟踪问题，本章的策略学习对象是围绕同一负载目标展开的联合协同行为，负载任务状态、单机局部状态、联合受力组织以及终端运输效果在强化学习建模中保持统一，状态、观测、动作和奖励的定义也围绕这一主线展开。为在强化学习框架下统一描述协同运输任务，每架无人机被视为具有独立动作输出的智能体，但其可直接获取的信息仅覆盖本机状态、负载相对状态与局部构型；负载运输效果则由三机联合动作共同决定。该任务可写成共享回报的 Dec-POMDP：
\begin{equation}
    \mathcal G=\langle \mathcal I,\mathcal S,\{\mathcal O_i\}_{i\in\mathcal I},
    \{\mathcal A_i\}_{i\in\mathcal I},P,r,\gamma\rangle,
\end{equation}
其中，$\mathcal I=\{1,2,3\}$ 为智能体集合，$\mathcal S$ 为全局状态空间，$\mathcal O_i$ 为第 $i$ 个智能体的局部观测空间，$\mathcal A_i$ 为其动作空间，$P$ 为状态转移函数，$r$ 为共享奖励函数，$\gamma$ 为折扣因子。

在该 Dec-POMDP 中，系统状态转移由联合动作共同驱动，奖励按照负载整体运输效果统一给出。时刻 $t$ 的全局状态写为
\begin{equation}
    s_t = \Phi\big(p_L,v_L,p_g,\{p_i,v_i,R_i,\Omega_i\}_{i=1}^{3}\big),
\end{equation}
其中 $\Phi(\cdot)$ 表示由物理状态到强化学习状态特征的映射。第 $i$ 个智能体在时刻 $t$ 的局部观测记为 $o_i^t$，联合动作记为
\begin{equation}
    a_t = (a_1^t,a_2^t,a_3^t).
\end{equation}
环境在执行联合动作后完成状态转移，并向全部智能体返回同一即时奖励 $r_t$。

执行阶段的单机观测属于局部可观测信息，策略面对的是带有历史依赖的局部决策问题。记第 $i$ 个智能体在时刻 $t$ 的局部信息历史为
\begin{equation}
    h_i^t = (o_i^0,a_i^0,\dots,o_i^{t-1},a_i^{t-1},o_i^t),
\end{equation}
则去中心化策略写为
\begin{equation}
    a_i^t \sim \pi_i(\cdot|h_i^t).
\end{equation}
实现中采用有限历史观测堆叠近似表征该历史依赖关系，训练阶段则通过全局状态价值评估显式刻画多机耦合，形成集中训练、去中心化执行的任务结构。
该 Dec-POMDP 表述与协同运输的物理对象保持一致。负载位置与速度、无人机状态以及目标位置共同组成全局状态；每架无人机依据自身与负载的相对关系、本机姿态和邻机相对位置构造局部观测；联合动作再经绳索耦合动力学共同作用于负载运输过程，训练阶段的全局价值评估与执行阶段的局部策略决策在这一框架内获得统一表达。在上述任务表述基础上，局部观测既要反映运输任务进展，又要保留绳索耦合与多机构型信息。仅输入本机状态不足以判断负载是否朝目标区域收敛，直接使用全局状态又不符合部署阶段的信息条件。基于协同运输的局部可得量，第 $i$ 架无人机的局部观测设计为
\begin{equation}
    o_i =
    \big(
    p_g-p_L,\;
    v_L,\;
    p_i-p_L,\;
    v_i,\;
    \bar q_i,\;
    \Omega_i,\;
    \{p_j-p_i\}_{j\ne i}
    \big),
\end{equation}
其中，$\bar q_i$ 表示无人机姿态的四元数表示，$\{p_j-p_i\}_{j\ne i}$ 表示其他无人机相对于本机的相对位置集合。
该观测向量将负载相对目标状态、本机相对负载状态和邻机相对位形统一写入单一输入。目标相对位置 $p_g-p_L$ 与负载速度 $v_L$ 描述运输进展，$p_i-p_L$、$v_i$、$\bar q_i$ 与 $\Omega_i$ 表示本机与负载之间的耦合状态，$\{p_j-p_i\}_{j\ne i}$ 则对应多机构型信息。为保留瞬态过程中的短时动态特征，局部观测在时间维上引入有限历史窗口，固定顺序拼接后的观测向量在批量训练中保持一致张量形状，策略可借助短时状态变化表征摆动趋势与收敛趋势；环境侧对异常值进行截断和清洗，训练阶段配合统计归一化处理，以削弱不同量纲特征对优化稳定性的影响。在该观测设计下，运输目标、本机耦合状态和多机构型约束均能够在局部信息条件下得到表达，去中心化执行结构得以保留，多机协同行为所需的关键信息也能够在单机决策中持续出现。在观测空间确定后，策略输出进一步定义为任务层控制参考，而非电机转速。多机协同受力与负载运动调节由策略层给出，高频姿态稳定与控制分配则交由确定性控制器完成。第 $i$ 架无人机的强化学习动作定义为 $a_i=(a_i^f,a_i^\psi)$，其中 $a_i^f=(a_{x,i},a_{y,i},a_{z,i})$ 表示三维归一化合力指令，$a_i^\psi$ 表示归一化偏航指令。策略输出先进行饱和限制，记 $\bar a_i^f = \mathrm{clip}(a_i^f,-1,1)$、$\bar a_i^\psi = \mathrm{clip}(a_i^\psi,-1,1)$，再映射为世界坐标系下的期望合力向量 $f_i^c = f_{i,\mathrm{mid}} + f_{i,\mathrm{half}}\odot \bar a_i^f$ 和期望偏航角 $\psi_i^c = \psi_{i,\mathrm{mid}} + \psi_{i,\mathrm{half}}\,\bar a_i^\psi$。其中，$f_{i,\mathrm{mid}}$ 与 $f_{i,\mathrm{half}}$ 分别表示动作映射的中心量与半量程，$\psi_{i,\mathrm{mid}}$ 与 $\psi_{i,\mathrm{half}}$ 分别表示偏航指令的中心量与半量程。

竖直方向中心量取为协同悬停状态下单架无人机承担的名义合力，记为 $\bar f_i = g\big(m_i+\frac{m_L}{3}\big)$。竖直控制区间围绕该中心量展开，水平合力上界由允许倾角和总推力能力共同确定。若记指令姿态上界为 $\vartheta_c$，则水平合力幅值上界写为 $f_{xy,i}^{\max} = \min\big(\bar f_i\tan\vartheta_c,\; f_i^{\max}\sin\vartheta_c\big)$。动作空间与四旋翼可实现工作区间保持一致。

期望合力向量 $f_i^c$ 定义在世界坐标系中，由其方向与偏航参考共同构造期望姿态。记
\begin{equation}
    \begin{aligned}
        b_{3,i}^d &= \frac{f_i^c}{\|f_i^c\|}, \qquad
        h_i = (\cos\psi_i^c,\sin\psi_i^c,0), \\
        b_{2,i}^d &=
        \frac{b_{3,i}^d\times h_i}{\|b_{3,i}^d\times h_i\|}, \qquad
        b_{1,i}^d = b_{2,i}^d\times b_{3,i}^d, \\
        R_i^d &= [\,b_{1,i}^d,\; b_{2,i}^d,\; b_{3,i}^d\,].
    \end{aligned}
\end{equation}
SE(3) 控制器根据当前姿态矩阵 $R_i$ 与期望姿态矩阵 $R_i^d$ 构造姿态误差 $e_{R,i} = \frac{1}{2}\big(R_i^{d\top}R_i - R_i^\top R_i^d\big)^\vee$，并结合当前角速度得到力矩参考 $\tau_i^c = -K_R e_{R,i} - K_{\Omega}\Omega_i$。总推力参考取期望合力模长并投影到允许集合，即 $f_i = \mathrm{clip}\big(\|f_i^c\|,\mathcal F_i\big)$，第 $i$ 架无人机的广义力参考写为 $w_i^c = (f_i,\tau_i^c)$。
广义力参考经控制分配矩阵映射到四个旋翼推力，记 $u_i^c = A_i^{-1} w_i^c$，再投影到执行器约束集合得到实际推力命令 $u_i = \Pi_{\mathcal U_i}(u_i^c)$。根据旋翼升力模型，电机角速度满足 $\omega_i = \sqrt{u_i/c_f}$，其中 $c_f$ 为推力系数。策略层合力意图、姿态控制与执行器约束在该分层链路中保持明确对应关系。
世界坐标系下的期望合力直接对应负载推进、升沉补偿和构型修正，偏航通道则约束机体朝向变化。动作向量中的水平分量调节推进方向与多机构型，竖直分量在名义悬停分担基础上补偿负载升沉与绳索张力变化。该动作空间保持了协同控制所需的表达能力，同时将策略输出限制在物理含义明确的受力空间内，电机级控制量由几何控制器和控制分配矩阵解析生成，策略优化由此集中于负载运输过程中的联合受力组织。

在状态、观测与动作均完成定义后，回报设计还需同时刻画运输过程中的持续改进与回合层面的成败结果。为使策略在长回合中获得连续优化信号，并在任务完成与失败边界处得到明确反馈，采用共享奖励机制，总奖励写为
\begin{equation}
    r_t = r_t^{\mathrm{step}} + r_t^{\mathrm{term}} - c_t.
\end{equation}
其中，$r_t^{\mathrm{step}}$ 为步进奖励，$r_t^{\mathrm{term}}$ 为成败最终奖励，$c_t$ 为代价项。步进奖励写为
\begin{equation}
    r_t^{\mathrm{step}} =
    r_t^{\mathrm{PBRS}}
    + r_t^{\mathrm{prog}}
    + r_t^{\mathrm{track}}
    + r_t^{\mathrm{refine}}.
\end{equation}
其中，势差奖励写为
\begin{equation}
    r_t^{\mathrm{PBRS}} = \Phi(s_{t+1})-\Phi(s_t),
\end{equation}
势函数定义为
\begin{equation}
    \Phi(s_t) = -\alpha_p \rho(\|e_p\|)-\alpha_v \rho(\|e_v\|),
\end{equation}
其中 $\rho(x)=\sqrt{x^2+\varepsilon_\rho^2}-\varepsilon_\rho$。
势函数由位置误差项与速度误差项加权构成。位置误差对应负载到目标区域的接近程度，速度误差刻画目标附近的穿越与振荡趋势；误差减小时势函数单调增大，势差奖励也随之与运输目标保持同向。平滑函数 $\rho(x)$ 在大误差区间近似绝对值增长，在零附近保持可导，运输推进阶段与精细调整阶段都能够获得稳定的梯度信号。
位置进展奖励写为
\begin{equation}
    r_t^{\mathrm{prog}}
    =
    \alpha_{\mathrm{prog}}
    \,\mathrm{clip}\big(\|e_p^{t}\|-\|e_p^{t+1}\|,-\bar c_p,\bar c_p\big),
\end{equation}
稳定运输奖励写为
\begin{equation}
    r_t^{\mathrm{track}}
    =
    \alpha_1 \exp(-\beta_1\|e_p\|)
    + \eta(e_p)\Big(
    \alpha_2 \exp(-\beta_2\|e_v\|)
    + \alpha_3 \exp(-\beta_3\theta_{\max})
    \Big),
\end{equation}
其中位置门控函数满足 $\eta(e_p)=\exp(-\kappa_p\|e_p\|)$。精细收敛奖励写为
\begin{equation}
    r_t^{\mathrm{refine}}
    =
    \alpha_r
    \Big[
    \log(\|e_p^{t}\|+\varepsilon)
    -
    \log(\|e_p^{t+1}\|+\varepsilon)
    \Big].
\end{equation}
步进奖励在回合内持续刻画运输过程。势差奖励与位置进展奖励推动负载朝目标区域收敛，稳定运输奖励约束目标附近的速度与摆动，精细收敛奖励进一步强化末端调整。位置门控函数使速度收敛和摆动抑制项在负载接近目标后占据更高权重，奖励重心也随运输阶段自然转向稳定保持。
成败最终奖励写为 $r_t^{\mathrm{term}} = r_t^{\mathrm{succ}} + r_t^{\mathrm{fail}}$。其中，$r_t^{\mathrm{succ}}$ 在负载位置误差、速度误差和飞行稳定性满足成功判据时触发，$r_t^{\mathrm{fail}}$ 在碰撞风险、高度越界、绳索摆角越界、姿态失稳或数值异常等情形下触发。成功奖励对应“进入目标区域并稳定保持”的完整任务闭环，失败奖励则把碰撞、失稳和严重摆动等不可接受行为直接标记为负例。
代价项写为
\begin{equation}
    c_t =
    c_t^{\omega}
    + c_t^{\mathrm{att}}
    + c_t^{\mathrm{eng}}
    + c_t^{\Delta u}
    + c_t^{\mathrm{time}}.
\end{equation}
其中，$c_t^{\omega}$ 对机体角速度进行惩罚，$c_t^{\mathrm{att}}$ 对姿态越界趋势进行惩罚，$c_t^{\mathrm{eng}}$ 对控制输入幅值进行惩罚，$c_t^{\Delta u}$ 对相邻时刻控制变化进行惩罚，$c_t^{\mathrm{time}}$ 为时间代价。动作变化代价写为 $c_t^{\Delta u} = \alpha_f \|\Delta f_t\|^2 + \alpha_{\tau}\|\Delta \tau_t\|^2$。
角速度、姿态越界趋势、控制能量、动作变化率和回合时长均被纳入代价项，运输过程中的激进机动据此受到抑制。高频控制变化和大姿态动作虽然可能带来短时位置改善，但在绳索耦合条件下容易放大负载摆动并降低长期稳定性。共享回报机制下，奖励与代价在每一时刻以统一标量形式反馈给全部智能体，局部决策始终围绕同一运输目标展开。

\section{去中心化多智能体强化学习算法}

\subsection{算法选型与适用性分析}

协同运输策略优化同时具有连续动作、共享回报和局部决策三项特征。将三架无人机及负载整体并为单一智能体后，联合动作维度会随系统规模迅速增加，策略搜索与部署复杂度随之上升；完全独立训练又难以稳定刻画绳索耦合、张力重分配与共享目标。采用 MAPPO，在训练阶段利用全局状态构造集中式价值评估，在执行阶段保留单机局部策略。MAPPO 的训练结构与协同运输任务的物理特征保持一致，集中式 critic 使用全局状态表征负载误差、多机构型和绳索耦合关系，actor 仍以局部观测生成单机动作，局部执行约束和系统级协同评价在同一优化框架内获得统一表达，共享回报与联合交互行为也能够在训练阶段稳定进入价值评估过程。

\subsection{策略网络与价值网络结构}

三架无人机在动力学形式、控制量表示和观测构成上具有同构特征，策略网络采用参数共享结构。设共享策略参数为 $\theta$，则第 $i$ 个智能体的策略记为 $\pi_{\theta}(a_i|o_i)$。共享 actor 接收带有限历史窗口的局部观测，输出归一化合力指令和归一化偏航指令的连续动作分布参数，参数共享使不同无人机在统一网络中复用样本，网络表征重点也集中于负载相对状态、绳索耦合与局部构型差异。价值网络采用集中式结构，其输入为全局状态 $s_t$，输出为状态价值 $V_{\phi}(s_t)$，其中 $\phi$ 为 critic 参数，负载误差、绳索摆动趋势与多机构型变化在价值评估阶段统一进入 critic，为优势估计和策略更新提供稳定基线。在该框架中，actor 依据局部观测生成单机动作，critic 依据系统整体状态评估长期运输收益，参数共享 actor 对多机之间的共性控制规律进行统一表征，集中式 critic 则持续吸收系统层面的耦合信息，两者共同对应协同运输中的局部执行与全局评价结构。

\subsection{优势估计与策略更新机制}

MAPPO 的更新建立在集中式价值评估与 PPO 裁剪优化之上。协同运输中的动作效应沿时间传播，某一时刻的推力调整需经绳索张力变化、负载速度演化与摆角传播后，才会体现为位置误差改善。对共享回报任务，时刻 $t$ 的时序差分残差写为
\begin{equation}
    \delta_t = r_t + \gamma V_{\phi}(s_{t+1}) - V_{\phi}(s_t).
\end{equation}
广义优势估计将不同时间尺度上的残差信息沿时间递推累积，得到
\begin{equation}
    \hat A_t = \sum_{l=0}^{T-t-1}(\gamma\lambda)^l\delta_{t+l},
\end{equation}
相应的回报目标为
\begin{equation}
    \hat R_t = \hat A_t + V_{\phi}(s_t).
\end{equation}
对第 $i$ 个智能体，裁剪策略目标写为
\begin{equation}
    L_i^{\mathrm{clip}}(\theta)
    =
    \mathbb E_t
    \Big[
    \min\Big(
    r_t^i(\theta)\hat A_t,\;
    \mathrm{clip}\big(r_t^i(\theta),1-\epsilon,1+\epsilon\big)\hat A_t
    \Big)
    \Big],
\end{equation}
其中，
\begin{equation}
    r_t^i(\theta)
    =
    \frac{\pi_{\theta}(a_i^t|o_i^t)}
    {\pi_{\theta_{\mathrm{old}}}(a_i^t|o_i^t)}.
\end{equation}
参数共享条件下，actor 总体损失写为
\begin{equation}
    L_{\mathrm{actor}}(\theta)
    =
    -\frac{1}{|\mathcal I|}
    \sum_{i\in\mathcal I}L_i^{\mathrm{clip}}(\theta)
    -\beta \mathcal H(\pi_{\theta}),
\end{equation}
critic 损失写为
\begin{equation}
    L_{\mathrm{critic}}(\phi)
    =
    \mathbb E_t\Big[
    \big(V_{\phi}(s_t)-\hat R_t\big)^2
    \Big].
\end{equation}
重要性采样比率度量新旧策略在同一历史样本上的动作概率变化幅度，裁剪操作将更新限制在有限邻域内，负载轨迹和绳索摆动对大步长更新的敏感性据此得到抑制。熵正则项维持探索性，critic 损失提高全局状态价值估计精度，共享回报、多步优势估计与裁剪更新共同形成协同运输任务中的多智能体优化过程。该更新机制与协同运输的时域耦合特征保持一致，单次推力调整往往在若干时刻后才转化为负载位置误差的改善，长期收益估计在策略优化中具有关键作用；已经形成的协同受力模式又对更新步长较为敏感，PPO 裁剪机制能够维持训练过程中的稳定性。

\subsection{训练流程}

MAPPO 的训练建立在批量并行环境之上，采样和更新均以环境批维为基本组织单位。环境经批量组织、定长回合组织和自动重置机制处理后，可同时推进 $B$ 个并行样本。设单次采样轨迹长度为 $H$，智能体数为 $N$。在第 $k$ 个采样时刻，环境返回批量局部观测张量 $O_k \in \mathbb R^{B\times N\times d_o}$ 以及对应的全局状态张量 $S_k \in \mathbb R^{B\times d_s}$。共享 actor 在批维和智能体维上同步推理，得到动作张量 $A_k \in \mathbb R^{B\times N\times d_a}$；局部动作沿智能体维拼接为联合动作张量 $A_k^{\mathrm{joint}} \in \mathbb R^{B\times (N d_a)}$，并输入环境完成一步批量状态转移。环境输出共享奖励 $R_k \in \mathbb R^{B}$、终止标志 $D_k \in \mathbb R^{B}$ 以及下一时刻状态。终止样本在批维上自动替换为新的初始状态，单次采样阶段获得的数据整理为
\begin{equation}
    \mathcal D=
    \{O_{0:H-1},S_{0:H},A_{0:H-1},R_{0:H-1},D_{0:H-1},\log \Pi_{0:H-1}\},
\end{equation}
其中 $\log \Pi_{0:H-1}$ 表示各时刻局部动作对应的对数概率信息。
长度为 $H$ 的采样轨迹生成后，采样缓存同时保留时间维、批维和智能体维。集中式 critic 在时间维和批维上计算状态价值，再结合终止标志沿时间维递推优势估计和回报目标，共享优势也被同步分配到对应的智能体样本上。更新阶段，局部观测、动作、旧策略对数概率和优势被整理为 actor 样本，全局状态与回报目标被整理为 critic 样本，并划分为若干小批量在同一轮采样数据上重复更新。时间维、批维和智能体维在训练流程中保持统一数据结构，联合行为的时间相关性能够在采样缓存中被完整保留，优势估计、价值评估与小批量重组均围绕同一数据布局展开，批量自动重置机制吸收了不同回合长度带来的差异，采样缓存始终维持一致张量形状。

\section{并行训练加速}

\subsection{并行训练需求分析}

协同运输环境的单步状态转移同时涉及三架无人机、一个负载和三根绳索的耦合动力学。负载位置调节、无人机姿态控制、动作映射、绳索约束和终止判定均在同一步中完成，单环境样本采集成本明显高于单机飞行控制问题。强化学习训练又依赖大规模交互样本估计策略梯度，环境吞吐量直接影响训练周期与样本分布覆盖范围，高吞吐量并行环境因而构成训练加速实现的基础条件。协同运输任务中的样本需求同时体现为数量规模与状态分布广度，策略需要覆盖不同初始构型、不同目标位置和不同瞬态摆动状态，负载位置误差、绳索摆角与多机相对构型之间的耦合关系才能在训练中被充分呈现；并行环境不仅缩短训练时间，也扩大了训练分布并提高样本利用效率。相较于单机悬停或路径跟踪问题，协同运输更加依赖状态覆盖，不同初始构型下的最优受力分配并不相同，训练若长期停留在少量场景上，策略容易形成依赖特定初始条件的局部控制模式，并行训练环境通过大量随机初始化回合同步展开，将推进、减速、摆动抑制与稳定保持等阶段性行为纳入更宽的状态分布。

\subsection{MJX 的并行计算机制}

MJX 为并行环境实现提供了统一的计算组织方式。环境推进、策略推理、奖励计算和优势估计均可写成由数组变换组成的函数链，物理状态、观测缓存、奖励辅助量和终止标志均以显式张量保存，批量环境之间的差异仅体现为状态取值和随机初始化，状态更新规律保持一致，时间维采样轨迹与批维并行样本也能够纳入同一数值流程。MJX 将 MuJoCo 的刚体动力学、执行器模型和约束求解转换为适合批量求值的物理状态表示，四旋翼姿态演化、负载平动响应与绳索约束传力在同一物理更新过程中完成，环境步进在实现上表现为同一动力学映射对整批样本的并行求值；对本章任务而言，物理引擎状态、任务层状态和学习器输入被统一放入同一计算空间。环境与训练沿同一数据结构组织后，物理层状态无需在面向对象表示与训练张量之间反复转换，奖励计算、终止判定和历史观测更新也在同一数组结构内完成，多机、多绳索与负载耦合环境中的实现一致性由此得到保证。

\subsection{协同运输环境的 MJX 实现}

协同运输环境在 MJX 中被组织为显式状态驱动的批量环境。若记物理状态为 $x_t$，观测历史缓存为 $h_t^{\mathrm{obs}}$，前一时刻势函数值和位置误差分别为 $\Phi_{t-1}$ 与 $e_{p,t-1}$，成功保持计数与精细收敛阶段计数分别为 $n_t^{\mathrm{succ}}$ 与 $n_t^{\mathrm{ref}}$，上一时刻控制缓存记为 $u_{t-1}$，则扩展环境状态写为
\begin{equation}
    \Xi_t=
    \big(
    x_t,\,
    h_t^{\mathrm{obs}},\,
    \Phi_{t-1},\,
    e_{p,t-1},\,
    n_t^{\mathrm{succ}},\,
    n_t^{\mathrm{ref}},\,
    u_{t-1}
    \big).
\end{equation}
其中 $x_t$ 对应物理系统的位姿、速度和控制输入等动力学量，其余状态量对应奖励计算、终止判定和历史观测更新。批量并行环境中的每个样本均沿该统一结构推进。
复位阶段在批维上同步执行。每个并行环境依据独立随机种子生成负载初始位置、目标位置和三机初始构型，并初始化无人机姿态、线速度和角速度；观测历史缓存由初始观测复制得到固定长度时间窗口。推进阶段接收批量联合动作后，先完成动作到控制量的映射，再执行单步物理求解，并从更新后的物理状态中提取负载位置、负载速度、目标相对位置、无人机位置、无人机速度、无人机姿态和角速度等任务量。局部观测与全局状态随之更新，共享奖励、成功判据、失败判据和回合长度约束也在同一步中同步更新，终止样本在批维上立即替换为新的初始状态。奖励辅助量、历史观测缓存与上一时刻控制量被纳入扩展环境状态后，任务层逻辑与物理状态能够同步推进，势函数差分、精细收敛阶段判定、成功保持计数和动作变化代价均依赖前一时刻信息，将其统一写入状态向量后，环境步进函数即可在单步更新中完成物理演化与任务更新。

\subsection{绳索约束的并行化实现}

绳索约束的并行化实现延续了前述物理建模思路。批量环境推进时，绳索长度、绳索方向和约束力与无人机和负载状态一同进入单步动力学更新，绳索松弛与拉紧切换也在同一求解过程中完成；任务层使用的绳索信息则表现为几何量，绳索方向、绳索长度和最大摆角均由系统状态直接导出，并进入观测构造、奖励计算和终止判定。动力学层中的绳索约束决定系统演化方式，任务层中的绳索几何量则刻画运输稳定性与安全性，绳索建模、环境状态表示和强化学习任务定义在该实现中保持一致，物理求解与任务评价之间不存在分离的状态描述。绳索长度、方向、伸长速度与张力在批维和智能体维上统一向量化计算，并随动力学更新传播到负载和平动状态，绳索只承受拉力、拉紧与松弛切换等关键物理特征在批量实现中得到保留，不同并行样本也无需维护独立的约束求解分支，批量物理更新与批量任务评价始终围绕同一组绳索几何量展开。


\section{仿真结果}

前述建模与算法分析给出了协同运输任务的控制框架。本节在三机协同运输环境中验证动力学建模、控制量表示与并行训练框架的有效性。实验设置遵循前文给出的 MAPPO 训练与执行框架，其奖励结构、控制量表示和物理环境与前文定义保持一致，评估围绕负载位置调节、速度衰减、摆动抑制和多机构型保持展开。为保证结果分析的一致性，协同运输任务的数值实验均在 MJX 并行环境中完成，物理仿真采用 MuJoCo 刚体动力学求解，三架四旋翼无人机通过三根柔性绳索连接同一负载，策略输出归一化期望控制量，经 SE(3) 姿态控制律和控制分配关系转化为各旋翼推力命令，训练阶段采用批量并行环境组织样本，并在 MAPPO 框架下基于 PPO 裁剪目标进行策略优化；除特别标注外，本节结果均在同一组实验平台与参数配置下获得。表~\ref{tab:triquad_platform} 给出软硬件平台配置，表~\ref{tab:triquad_sim_setup} 汇总仿真模型、动力学参数与边界判据，表~\ref{tab:triquad_train_setup} 汇总训练超参数，后续训练收敛性分析、典型控制效果与扩展任务验证均在该组条件下进行。

\begin{table}[htbp]
    \centering
    \caption{三机协同负载运输实验平台配置}
    \label{tab:triquad_platform}
    \renewcommand{\arraystretch}{1.15}
    \begin{tabular}{p{0.26\textwidth}p{0.36\textwidth}p{0.18\textwidth}}
        \toprule
        参数 & 数值/规格 & 单位/说明 \\
        \midrule
        操作系统 & Ubuntu 22.04.5 LTS & 版本\\
        中央处理器 & Intel Core i9-13900KF & 型号 \\
        图形处理器 & NVIDIA GeForce RTX 4090 & 型号 \\
        显存容量 & 24 & GB\\
        主内存容量 & 62 & GB\\
        \bottomrule
    \end{tabular}
\end{table}

\begin{table}[htbp]
    \centering
    \caption{三机协同负载运输仿真实验参数}
    \label{tab:triquad_sim_setup}
    \renewcommand{\arraystretch}{1.15}
    \begin{tabular}{p{0.14\textwidth}p{0.42\textwidth}p{0.16\textwidth}p{0.16\textwidth}}
        \toprule
        符号 & 含义 & 数值 & 单位/说明 \\
        \midrule
        $T_c$ & 环境控制周期 & 0.02 & s \\
        $T_s$ & MuJoCo 物理积分步长 & 0.004 & s \\
        $N_{\mathrm{epi}}$ & 单回合仿真步数 & 1000 & 步 \\
        $m_{\mathrm{u}}$ & 单架无人机质量 & 0.9 & kg \\
        $J_{\mathrm{u}}$ & 单架无人机转动惯量对角元 & (0.02, 0.02, 0.04) & kg$\cdot$m$^2$ \\
        $m_{\mathrm{p}}$ & 负载质量 & 0.6 & kg \\
        $J_{\mathrm{p}}$ & 负载转动惯量对角元 & (0.01, 0.01, 0.01) & kg$\cdot$m$^2$ \\
        $l_{\mathrm{a}}$ & 四旋翼机臂长度 & 0.22 & m \\
        $L_0$ & 绳索弹簧原长 & 0.82 & m \\
        $k_{\mathrm{r}}$ & 绳索等效刚度 & 30 & N/m \\
        $c_{\mathrm{r}}$ & 绳索等效阻尼 & 2.0 & N$\cdot$s/m \\
        $f_{\max}$ & 单个旋翼最大推力 & 10 & N \\
        $F_{\min}$ & 集体推力下界 & 0 & N \\
        $F_{\max}$ & 集体推力上界 & 20 & N \\
        $\delta_p$ & 成功判据位置误差阈值 & 0.18 & m \\
        $\delta_v$ & 成功判据速度误差阈值 & 0.25 & m/s \\
        $\delta_{\omega}$ & 成功判据角速度阈值 & 1.0 & rad/s \\
        $T_{\mathrm{h}}$ & 成功保持时间阈值 & 0.5 & s \\
        $d_{\mathrm{coll}}$ & 无人机碰撞限制距离 & 0.44 & m\\
        $\theta_{\mathrm{r,max}}$ & 绳索摆角终止阈值 & 50 & 度 \\
        $\theta_{\mathrm{u,max}}$ & 机体姿态终止阈值 & 60 & 度 \\
        \bottomrule
    \end{tabular}
\end{table}

\begin{table}[htbp]
    \centering
    \caption{三机协同负载运输 PPO 训练超参数}
    \label{tab:triquad_train_setup}
    \renewcommand{\arraystretch}{1.15}
    \begin{tabular}{p{0.14\textwidth}p{0.42\textwidth}p{0.16\textwidth}p{0.16\textwidth}}
        \toprule
        符号 & 含义 & 数值 & 单位/说明 \\
        \midrule
        $N_{\mathrm{train}}$ & 总训练环境步数 & $2\times 10^8$ & 步 \\
        $N_{\mathrm{env}}$ & 并行训练环境数量 & 8192 & 个 \\
        $N_{\mathrm{rep}}$ & 动作重复次数 & 1 & 次 \\
        $H_{\mathrm{obs}}$ & 历史观测帧数 & 3 & 帧 \\
        $L_{\mathrm{u}}$ & 单次采样轨迹长度 & 40 & 步 \\
        $B$ & 每个小批量样本数 & 2048 & 个 \\
        $N_{\mathrm{mb}}$ & 每轮更新的小批量个数 & 32 & 个 \\
        $N_{\mathrm{upd}}$ & 每批数据的优化轮数 & 8 & 轮 \\
        $\alpha$ & Adam 学习率 & $2\times 10^{-4}$ & 无量纲 \\
        $\gamma$ & 折扣因子 & 0.997 & 无量纲 \\
        $\beta_{\mathrm{ent}}$ & 熵正则系数 & $1.0\times 10^{-3}$ & 无量纲 \\
        $s_r$ & 奖励缩放系数 & 5.0 & 无量纲 \\
        \bottomrule
    \end{tabular}
\end{table}

为量化训练过程与最终策略性能，本节采用成功率、平均奖励、平均到达时间和终端位置误差作为主要评价指标。设评估阶段共执行 $N_{\mathrm{eval}}$ 个回合，其中成功回合数记为 $N_{\mathrm{succ}}$，则成功率定义为
\begin{equation}
    P_{\mathrm{succ}} = \frac{N_{\mathrm{succ}}}{N_{\mathrm{eval}}}.
\end{equation}
平均奖励用于衡量策略在整个回合上的综合收益，记为
\begin{equation}
    \bar R = \frac{1}{N_{\mathrm{eval}}}\sum_{k=1}^{N_{\mathrm{eval}}} R_k,
\end{equation}
其中 $R_k$ 表示第 $k$ 个回合的累计回报。平均到达时间仅在成功回合上统计，记为
\begin{equation}
    \bar T_{\mathrm{arr}}=
    \frac{1}{N_{\mathrm{succ}}}
    \sum_{k\in\mathcal S} T_k,
\end{equation}
其中 $\mathcal S$ 为成功回合集合，$T_k$ 表示第 $k$ 个成功回合达到目标区域并满足稳定保持判据所需时间。终端位置误差定义为各评估回合终止时刻的负载位置误差平均值，即
\begin{equation}
    \bar e_p^{\mathrm{term}}=
    \frac{1}{N_{\mathrm{eval}}}
    \sum_{k=1}^{N_{\mathrm{eval}}}
    e_{p,k}^{\mathrm{term}}.
\end{equation}
四项指标分别对应任务完成概率、综合控制质量、运输效率与终端稳态精度，训练收敛性分析、典型回合时域响应和扩展场景验证均围绕这四项指标展开。在上述实验设置与评价指标基础上，进一步考察策略在优化过程中的性能演化。图~\ref{fig:triquad_train_curves} 给出了训练步数增加时成功率、平均奖励、平均到达时间和终端位置误差的变化。训练初期，策略尚未形成稳定的多机协同受力与摆动抑制规律，成功率和平均奖励处于较低水平；随着采样规模增加，策略逐步学会利用三架无人机的联合受力调整负载运动方向，成功率持续上升，平均奖励同步提高，奖励结构中的推进、稳定与终端判据也在训练中形成方向一致的优化信号。平均到达时间和终端位置误差反映运输效率与末端精度的变化，训练前期，部分回合虽能接近目标，但到达时间较长，目标附近振荡也较为明显；训练推进后，成功回合的平均到达时间逐步缩短，终端位置误差持续减小，策略能够在把负载送达目标附近的同时完成减速与稳定保持。四项指标在训练过程中呈现出较明确的阶段性特征，早期阶段，策略主要消除显著失稳和无效受力，成功率提升最为明显；中期阶段，平均奖励和平均到达时间改善更快，联合推进模式逐步稳定；后期阶段，终端位置误差继续下降，优化重点转向目标附近的减速与稳态保持，训练结果表明，并行环境中的 PPO 训练能够逐步形成稳定有效的协同运输规律。

\begin{figure}[htbp]
    \centering
    \includegraphics[width=1\linewidth]{pic/part3pic/tri_quad_payload_training_curves.png}
    \caption{训练过程中主要评价指标随训练步数的变化}
    \label{fig:triquad_train_curves}
\end{figure}

在训练过程表现出稳定收敛后，进一步选取最终策略在标称协同运输场景下执行的一次典型回合，对时域响应与空间轨迹进行联合分析。图~\ref{fig:triquad_time_response} 给出了相对位置误差、负载速度、最小无人机间距和最大绳索摆角的时间响应；图~\ref{fig:triquad_typical_traj} 给出了对应回合的三维轨迹图。时域响应表明，负载相对目标的位置误差在运输过程中持续下降，负载速度在中后期逐步衰减，控制过程经历了由有效推进到目标附近减速与保持的演化；最小无人机间距始终高于安全阈值，最大绳索摆角在瞬态调整过程中出现波动，但整体保持在允许范围内。三维轨迹图中，三架无人机围绕负载维持稳定运输构型，负载轨迹由初始位置平滑收敛至目标区域附近，机间轨迹未出现明显交叉或发散。典型回合中的时域响应与空间轨迹相互对应，策略在运输前段形成指向目标的有效合力，接近目标后逐步提高对速度与摆动的抑制强度；多架无人机在绳索耦合条件下形成稳定支撑关系，空间轨迹中的平滑收敛也与时域响应中的速度衰减和摆角受控保持一致，最终策略已经学到与绳索耦合动力学相适应的协同运输行为模式。表~\ref{tab:triquad_final_result} 汇总了最终训练轮次在评估阶段得到的成功率、平均奖励、平均到达时间和终端位置误差，作为主任务实验的最终定量结果。

\begin{figure}[htbp]
    \centering
    \includegraphics[width=1\linewidth]{pic/part3pic/tri_quad_payload_time_response.png}
    \caption{最终策略在典型回合下的时域响应}
    \label{fig:triquad_time_response}
\end{figure}

\begin{figure}[htbp]
    \centering
    \includegraphics[width=0.6\linewidth]{pic/part3pic/tri_quad_payload_trajectory_3d.png}
    \caption{最终策略在典型回合下的三维协同运输轨迹}
    \label{fig:triquad_typical_traj}
\end{figure}

\begin{table}[htbp]
    \centering
    \caption{最终训练轮次的主要评估结果}
    \label{tab:triquad_final_result}
    \renewcommand{\arraystretch}{1.15}
    \begin{tabular}{p{0.34\textwidth}p{0.18\textwidth}p{0.18\textwidth}}
        \toprule
        指标 & 结果 & 单位/说明 \\
        \midrule
        成功率 & 100 & \%\\
        平均奖励 & 54.85 & 无量纲\\
        平均到达时间 & 1.71 & s\\
        终端位置误差 & 0.005 & m\\
        \bottomrule
    \end{tabular}
\end{table}

表中的四项结果对应主任务实验的定量摘要。成功率与平均奖励概括策略的总体完成质量，平均到达时间反映运输效率，终端位置误差刻画目标附近的最终保持精度。

在完成定点协同运输任务验证后，保持策略参数不变，并将目标点在线更新为平面圆周轨迹，以检验最终策略对时变参考的响应能力。该实验不作为主任务性能评价依据，重点考察训练分布之外的协同表现。在圆周轨迹扩展验证中，目标点沿水平圆周连续变化，负载与三架无人机需要在维持协同构型的同时持续调整运输方向。图~\ref{fig:triquad_circle_response} 给出了时域响应，图~\ref{fig:triquad_circle_traj} 给出了三维轨迹。圆周参考下的负载位置误差和最大绳索摆角通常大于定点运输场景，目标持续变化时，系统需要更频繁地调整合力方向并修正姿态；三机对负载的基本协同支撑关系保持稳定，负载轨迹能够在一定范围内沿参考路径演化。扩展验证表明，定点运输任务中训练得到的策略已经形成一定的运输与构型保持规律，对时变参考仍能维持基本负载跟随与多机协同支撑关系；与此同时，圆周参考下的位置误差和摆角响应也反映出当前训练目标仍以定点到达和稳定保持为主，对复杂时变轨迹的适应能力仍然有限。进一步提高时变参考下的跟踪性能，还需要在训练任务表述中显式引入时变目标分布、轨迹跟踪奖励或更丰富的课程训练机制。圆周参考要求策略持续调整合力方向，目标方向在整个回合中连续旋转，三架无人机在维持负载支撑关系的同时，还需持续调整切向推进方向和法向约束关系；局部观测中的相对位置信息、控制量表示中的水平合力分量以及奖励中的速度与摆动抑制项均在该场景下持续发挥作用，扩展场景中的基本协同关系保持稳定，前述任务建模对静态目标之外的时变参考也表现出一定的结构迁移能力。

\begin{figure}[htbp]
    \centering
    \includegraphics[width=1\linewidth]{pic/part3pic/circle_time_response.png}
    \caption{复杂轨迹扩展验证中的时域响应}
    \label{fig:triquad_circle_response}
\end{figure}

\begin{figure}[htbp]
    \centering
    \includegraphics[width=0.6\linewidth]{pic/part3pic/circle_trajectory_gradient.png}
    \caption{复杂轨迹扩展验证中的三维协同轨迹}
    \label{fig:triquad_circle_traj}
\end{figure}

\section{本章小结}

本章围绕低空物流场景中的多无人机协同负载运输控制问题，建立了绳索约束下三机——负载系统的耦合动力学模型，并将控制过程写成共享回报的多智能体序贯决策问题。围绕该任务，给出了局部观测、归一化控制量表示与奖励函数，算法部分讨论了 MAPPO 形式的集中训练、去中心化执行框架，实验层面则基于参数共享 PPO 与 MJX 并行环境完成训练评估。仿真结果表明，所构建的协同运输控制框架能够在复杂耦合条件下实现负载到达、速度衰减、摆动抑制与构型保持，并行环境也为策略训练提供了高吞吐量样本支撑。定点运输结果与圆周参考扩展验证共同说明，该方法在多无人机协同负载运输任务中具有较好的稳定控制能力，并为后续面向复杂时变任务的扩展提供了基础。
